%-*-tex-*-
\ifundefined{writestatus} \input status \relax \fi %
\chcode{footnot}
\def\cqu{}

{\catcode`\?=7
\global\def\ftn{$\sphantom{\cheadfont S}?\ddagger$} 
}
\autonumberingon

\chapterhead{footnot}{FOOTNOTES\ftn}
\sfootnote{\ddagger}{This really should not have been footnoted. However to
 obtain the correct placement of the $\ddagger$, 
 |$\sphantom{\cheadfont S}^\ddagger$| was actually used. This footnote leads 
to a large number of problems. It will show up in the automatic listing of
the table of contents and the header text. To prevent this, the 
|\let\ftn=\relax| is placed immediately after the |\chapterhead| and 
just before |\maketoclist| produces the table of contents.
}
\let\ftn=\relax
\intex\ supplies a simple but powerful 
footnote\footnote{\dagger}{For making footnotes like this, even with
imbedded mathematics, {\it ie} $\int_0^{\infty} g(x)dx$, or {\bf font}
changes.}
capability. This replaces the |\footnote| command in {\it plain}. {\bf Do not
use |\footnote| in  |\title| or any of the section head commands.} 
The footnote will  get lost. If you must footnote a title or section
head, then use |\up{<mark>}| where you call out the footnote and |\sfootnote|
with the same |<mark>| in the vertical list of the surrounding {\bf page}.
For those of you who wish to change the default formats there is an 
|\everyfootnote={<token list>}| that will be inserted at the beginning of
every footnote list. 
\shead{footnotcomlist}{Command List}
\begintwocolumn
\pla|\dimen\footins|
\ext|\everyfootnote|
\ext|\footnote|
\ext|\footnotefont|
\pla|\footnoterule|
\pla|\skip\footins|
\ext|\sfootnote|
\ext|\up|
\endthreecolumn

\shead{footnotcomforms}{Command Forms}
\beginblockmode
\ext\@|\everyfootnote = {<token list>}|
\nbr
The |<token list>| is inserted before every footnote list. It could be used,
for instance, to change the spacing between lines or any other default
feature that is found undesirable. 

\mbr
\ext\@|\footnote{<math mode mark>}{<text>}|
\nbr
The |<math mode mark>| is the marker  that is placed
on the word that is being footnoted. Note that it is in {\bf math mode}. This 
command should be used only in (internal) horizontal mode. 
If the footnote gets lost for some reason, use the |\up{<math mode mark>}| at
the appropriate place in the text and |\sfootnote| with the same 
|<math mode mark>|. The |\sfootnote| should be in the vertical list of the
main page.
\mbr
\ext\@|\footnotefont|
\nbr
The footnote font (really an entire family) 
\dots it is set when the document font
is set or specifically. 
\mbr
\pla\@|\footnoterule|
\nbr
The footnote line at the bottom of a page. This is modified by redefining it
through a |\def\footnoterule{<new definition>}|.
\mbr
\ext\@|\up{<math mode mark>}|
This command  superscripts the |<math mode mark>| to the preceding character.
It assumes that the current font family is active. Since the
|\cheadfont| in the title is much bigger than the default |\tenpoint| of this
book, |{\cheadfont FOOTNOTES\up{\ddagger}}| would look like this
{\cheadfont FOOTNOTES\up{\ddagger}}.

\mbr
\pla\@|\skip\footins [=] <dimen>|
\nbr
This sets the initial skip before the first
footnote to |<dimen>|.
\mbr
\pla\@|\dimen\footins [=] <dimen>|
\nbr
This sets the maximum height allowed for
footnotes to |<dimen>|. If this maximum height is exceeded, the footnotes spill over to the
next page. It is set in \intex\  at 10in. 

\mbr
\ext\@|\sfootnote{<math mode mark>}{<text>}|
\nbr
This should be used for footnotes that appear in the vertical mode. It is
necessary if the original footnote to be called out is too deep inside
vertical boxes. An example of this is the |\title|, the various section and
chapter heads, and the |\captionbox| in a |\figureinsert| or |\tableinsert|.
It should not be used in either horizontal or internal vertical mode.
\endblockmode
\ejectpage
\done